

% ============================================================
%  签证与居留许可（独立编译章节）
% ============================================================

\section{签证办理}

\subsection{阅读须知}

本章节信息仅为学联整理汇总，具体签证办理细请\textbf{务必}以丹麦移民局官网为准：

\begin{itemize}[label=•,leftmargin=2em,itemsep=2pt]
  \item 用于移民局网址：\href{https://www.nyidanmark.dk/en-GB/}{https://www.nyidanmark.dk/en-GB/}
  \item SIRI 联系方式汇总：\href{https://www.nyidanmark.dk/en-GB/Contact-us}{https://www.nyidanmark.dk/en-GB/Contact-us}


\end{itemize}


若想加入学联2025年丹麦留学新生群，获取一手信息，可在学联微信公众号【哥本哈根CSSA】或小红书【哥本哈根学联（CSSA）】后台私信，获取入群方式私信


\subsection{线上递签申请流程}

% -------- 线上递签申请流程（灰白条 + scriptsize）--------
\noindent
\scriptsize
\rowcolors{1}{gray!10}{white}           % ① 从第 3 行（数据区首行）开始灰白交替
\renewcommand{\arraystretch}{1.2}       % ② 行高微调，可按需调整
\begin{longtable}{p{0.24\linewidth}p{0.70\linewidth}}
  %\caption{线上递签申请流程}\label{tab:e-visa-steps}\\
  \toprule
  \rowcolor{gray!10} % <<< 加这一行，表头变灰
  \textbf{步骤} & \textbf{说明} \\ 
  \midrule
  \endfirsthead

  \multicolumn{2}{l}{\textit{接上页}}\\
  \toprule
  \textbf{步骤} & \textbf{说明} \\ 
  \midrule
  \endhead

  \midrule\multicolumn{2}{r}{\textit{续下页}}\\
  \bottomrule
  \endfoot

  \bottomrule
  \endlastfoot

  1. 提前准备文件 & 全本护照彩色扫描件（含封面、封底）；学校 offer；存款证明 / 学费缴纳证明（彩色扫描） \\
  2. 进入丹麦移民局签证 & \url{https://www.nyidanmark.dk/en-GB/Applying/Study} \\
  3. 选择签证类型 & 选择 \emph{Higher educational programme}（PhD选择\emph{PhD studies}），点击 \textbf{How to apply} \\
  4. 创建 case order ID & Create case order ID \\
  5. 付款（Case order ID 费用） & 约 2500 DKK，付款后务必打印 receipt \\
  6. 在线填写 ST1 表格 & 用学校邮件提供的账号密码登录；点击 \textit{"Use the online form ST1 – remember to save the screen receipt"} 完成填写，并打印 ST1 表格及回执 \\
  7. 上传所需文件 & Case ID receipt、学校 offer、学费缴费证明 / 存款证明等 \\
  8. 上传 Sworn Declaration & 系统生成宣誓证明，手写签名后扫描上传；关闭页面前再次确认已打印 ST1 表格、回执及 Case ID 收据 \\
  9. 缴纳签证费用 & \url{https://dys.um.dk/permit/} 支付约 1700 DKK；保留邮箱收到的 receipt 及 Terms \& conditions \\
  10. 14 天内采集生物信息 & 提交后 14 天内前往签证中心录指纹 / 拍照 \\
  11. 等待签证到达 & — \\

\end{longtable}
\normalsize


\paragraph{什么是 ST1 表格？}

ST1 是丹麦政府为学生居留许可设计的官方表格，分为两部分：\textbf{Part 2} 由学校填写，\textbf{Part 1} 由学生完成。仅在学生满足全部入学条件后，学校才会填写 ST1；应届生须拿到双证后尽快换 \emph{unconditional offer}，以便早日获取 ST1 并递签。

\begin{itemize}[label=•,leftmargin=2em,itemsep=2pt]
  \item 2024 年签证费：2490 DKK，建议使用 Visa 支付。
  \item 缴费后请务必截图／下载成功回执。
  \item 签证审核时长不定，建议学校发完 ST1 后尽快提交并录生物信息。
\end{itemize}

% --------------------------------------------------
\subsection{线下生物信息录入}

\subsubsection{线下递签材料准备 (From 2024)}

\tcbset{before skip=4pt,after skip=4pt}
\begin{itemize}[label=•,leftmargin=2em,itemsep=2pt]
  \item 护照原件 + \textbf{彩印}全本复印件 1 份（含空白页与封皮）+ 护照个人信息页单独 1 张
  \item 学校 con／uncon offer（建议拿到 uncon 后再递交）
  \item Tuition fee proof（学费缴费证明或存款证明）
  \item 1 张白底照片（3.5 cm × 4.5 cm）
  \item Case order ID 缴费收据
  \item 使馆签证申请费缴费收据
  \item Terms and conditions
  \item ST1 表格
  \item ST1 回执
  \item Sworn Declaration 宣誓证明
  \item 在校证明（仅 guest／exchange student 需要）
\end{itemize}

\begin{grayleftbox}{\textbf{学长学姐碎碎念}}
  \begin{itemize}
    \item 部分文件可使用无水印的手机扫描（如夸克扫描），扫描要注意确保打印出来的时候，所有页面的大小要符合实际材料的大小，否则也会要求现场重印。不过还是打印店更方便！

    \item 使馆要求 \textbf{所有材料彩印}，包括护照全本！含封面和封底。
  \end{itemize}
\end{grayleftbox}


\subsubsection{中国大陆签证中心地点}

\begin{table}[ht]
  \small
  \centering
  %\caption{中国大陆签证中心地点}\label{tab:visa-center-cn}
  \begin{tabularx}{\textwidth}{p{0.22\textwidth} X}
    \toprule
    \textbf{城市} & \textbf{地址} \\ \midrule
    北京   & 朝阳区光华路 9 号 \\
    上海   & 黄浦区四川中路 213 号 \\
    广州   & 天河区体育西路 189 号 城建大厦 3 楼 \\
    重庆   & 两江新区西湖支路 2 号 精信中心 B 塔 3 层 6 号 \\ 
    \bottomrule
  \end{tabularx}
\end{table}



\begin{grayleftbox}{\textbf{学长学姐碎碎念}}
\begin{itemize}
\item 上海下签还蛮快的，但办理时长因人而异，建议一拿到 con 就申请。
\item 如毕业证下发较晚，可请导师用学校邮箱出具可毕业证明先换 uncon（2021 年经验）。
\item 北京中心可 walk-in，约 30 分钟即可办完。
  \end{itemize}
\end{grayleftbox}


\section{居留许可 \& 黄卡申请}

新生入境后应尽快\textbf{线上预约 International House} 并\textbf{前往线下}办理居留证件，方可在丹麦正常学习、工作与生活。

\subsection{居留材料介绍}

\small
\begin{longtable}{p{0.22\linewidth}p{0.72\linewidth}}
\toprule
\rowcolor{gray!10} % <<< 加这一行，表头变灰
\textbf{材料} & \textbf{作用} \\
\midrule
\endhead
CPR number & 丹麦社会保障号，相当于身份证号码；无 CPR 无法开银行账户、就医、借书、纳税、领工资等 \\
黄卡 (sundhedskort) & 公共医疗保险卡 \& 身份证明，印有姓名、地址、CPR、家庭医生等信息；搬家需在 \href{https://lifeindenmark.borger.dk/housing-and-moving}{https://lifeindenmark.borger.dk/housing-and-moving} 更新地址并领取新黄卡 \\
粉卡 (opholdskort) & 居留卡 / ID card；非欧盟居民在丹麦合法居住的身份证明，含照片、签证类型及到期日；申根区旅行、丹麦出入境必备 \\
MitID & 丹麦国家数字身份认证系统，线上办事与支付的主要工具 \\
\bottomrule
\end{longtable}
\normalsize

\subsection{如何预约办理}

\small
\begin{longtable}{p{0.24\linewidth}p{0.70\linewidth}}
    \toprule
    \rowcolor{gray!10} % <<< 加这一行，表头变灰
    \textbf{步骤} & \textbf{说明} \\
    \midrule
    \endfirsthead
    \multicolumn{2}{l}{\textit{接上页}}\\
    \toprule
    \textbf{步骤} & \textbf{说明} \\
    \midrule
    \endhead
    \midrule\multicolumn{2}{r}{\textit{续下页}}\\\bottomrule
    \endfoot
    \bottomrule
    \endlastfoot

    1. 提前准备文件 &
      护照、居留许可 (resident permit letter)、住房合同 \\

    2. 线上申请 &
      \url{https://ihcph.kk.dk/registration-guidance/cpr-registration} \\

    3. 选择线下办理时间 &
      申请通过后，系统发送链接；选择合适时间并获取二维码／编码 \\

    4. 线下办理 &
      地址：International House Copenhagen, Nyropsgade 1, 1602 Copenhagen V \newline
      必带：护照、居留信、住房合同 \newline
      可带：学校 offer、存款证明/学费缴费证明
 \\

    5. 获得居留材料 &
      当场获得写有 CPR 号的小纸条；务必邮件告知学校 Student Service ；\newline
      黄卡、粉卡将在 2–4 周内分别邮寄至住址，请及时查收信箱 \\

    6. 激活 MitID &
      线下办理次日即可激活：下载 MitID App，按提示用护照完成激活 \\
\end{longtable}
\normalsize

% --- Comment Box -------------------------------------------------
\begin{grayleftbox}{\textbf{小贴士： }}
  \begin{itemize}
    \item 拿到 CPR 并激活 MitID 后，记得邮件告知学校，以便在学校系统里更新CPR；
    \item 搬家需及时 report 并更新黄卡地址、重新匹配家庭医生，逾期可能罚款。
    \item \textit{Digital Post} App 可查看政府来信；报税、交通罚款、青年卡政策变动等都在此处理。
    \item \textit{Sundhedskortet} App 可获得电子黄卡。
    \item 申根区旅行请随身携带护照 + 粉卡。
  \end{itemize}
\end{grayleftbox}
% ----------------------------------------------------------------





% ---------------- End of file ----------------